\section{Give a sensor record a portable representation}
The packet parser supplied a fixed layout, and the WAV reader discovered fields inside
chunks. This project gives a record a representation whose individual values describe
their own kind and extent. A receiver can find the next value without knowing its
machine address or the compiler layout of a structure. The application still needs a
schema: recognizing an integer does not tell the receiver whether it is a counter,
a length or a measurement.

Use CBOR to recover a bounded telemetry record, then write a deterministic encoding of
its required fields. The new work is an item-head reader, a bounded walker for nested
values, and a decoder that publishes a complete result only after validation. Keep
transport separate. The five entries in this chapter operate on memory; a real file
adapter can reuse the transfer work from the previous projects.

Begin with the fixed trace produced by
\path{drafts/cbor-fixtures/fixture.py}. Its input deliberately contains optional metadata
and some wider argument representations than the output writer will choose. From the public book root, run
\texttt{python3} \path{drafts/cbor-fixtures/fixture.py} \texttt{trace.cbor}.
Use a fresh filename with \texttt{--record} to inspect its 176-byte local record image.
The helper refuses an existing output. Compare
the recovered record before comparing serialized bytes. Equality of the required
application fields is the first observation; the writer's exact byte representation
is a second, separately specified observation.
